\documentclass[aspectratio=169]{beamer}

\usepackage[utf8]{inputenc}
\usepackage[T1]{fontenc}

\usepackage{pgfplots}
\usepackage[dvipsnames]{xcolor}
\usepackage{amssymb, amsmath, tabularx, booktabs}
\usepackage{amsfonts}
\usepackage{makecell}
\usepackage{tikz}
\usepackage{hyperref}

\usetikzlibrary{shadings}
\usetikzlibrary{positioning}
\usetikzlibrary{calc, arrows.meta}
\usetikzlibrary{shapes.misc}

\usepgfplotslibrary{colormaps, patchplots}

\pgfplotsset{compat = newest}

\usetheme{Tufte}

\DeclareMathOperator*{\argmax}{argmax}
\DeclareMathOperator{\sg}{sg}
\DeclareMathOperator{\symlog}{symlog}

\title{O que danado é um \textit{World Model}?}
\subtitle{Uma Introdução às Representações Internas em Inteligência Artificial}
\author{Vitor N. Cabral de Oliveira}
\institute{\href{mailto:vnco@cin.ufpe.br}{vnco@cin.ufpe.br}}
\date{\today}

\tikzset{
  box/.style={draw=wp-primary!50, fill=wp-bg, minimum width=1.2cm, minimum height=0.6cm, text=wp-text, align=center, font=\small},
  input/.style={draw=wp-muted!50, fill=wp-bg, minimum width=0.8cm, minimum height=0.4cm, text=wp-text, align=center, font=\small},
  loss/.style={draw=wp-accent!50, fill=wp-bg, minimum width=2.0cm, minimum height=0.6cm, text=wp-text, align=center, font=\small},
  emph/.style={draw=wp-primary!50, fill=wp-bg, minimum width=1.2cm, minimum height=0.6cm, text=wp-text, align=center, font=\small},
  layer/.style={draw=wp-primary!50, fill=wp-bg, minimum width=1.8cm, minimum height=0.6cm, text=wp-text, align=center},
  error/.style={draw=wp-accent!60, fill=wp-bg, minimum size=0.55cm, text=wp-text, font=\tiny},
  comp/.style={draw=wp-primary!50, fill=wp-bg, minimum width=2.0cm, minimum height=0.6cm, text=wp-text, align=center},
  era/.style={draw=wp-primary!50, fill=wp-bg, minimum width=2.4cm, minimum height=0.5cm, text=wp-text, align=center, font=\small}
}

\begin{document}

% ===== Title Frame =====
\begin{frame}[plain]
  \titlepage
\end{frame}

% ===== Outline =====
\begin{frame}{Roteiro}
  \tableofcontents
\end{frame}

% ============================================================
\section{Introdução}
% ============================================================

\begin{frame}{O que é um \textit{World Model}?}
  \begin{block}{Definição}
    Um \textbf{world model} é uma representação interna que um agente constrói do seu ambiente, permitindo:
  \end{block}
  \vspace{0.3em}
  \begin{itemize}
    \item \textbf{Prever} estados futuros do ambiente
    \item \textbf{Raciocinar} sobre consequências de ações
    \item \textbf{Planejar} sem tentativa-e-erro no mundo real
  \end{itemize}
  \vspace{0.5em}
  \begin{center}
    \begin{tikzpicture}[node distance=0.6cm and 1.2cm, auto, >=latex]
      \node[box, minimum width=2.2cm] (percep) {Percepção};
      \node[box, minimum width=2.2cm, right=of percep] (mem) {Memória};
      \node[box, minimum width=2.2cm, right=of mem] (plan) {Planejamento};
      \draw[->, wp-primary, thick] (percep) -- (mem);
      \draw[->, wp-primary, thick] (mem) -- (plan);
      \draw[->, wp-primary, thick, dashed] (plan.south) .. controls +(down:0.5cm) and +(down:0.5cm) .. (percep.south);
      \node[wp-muted, font=\tiny, below=1.2cm] at (mem.center) {loop percepção--ação};
    \end{tikzpicture}
  \end{center}
\end{frame}

% ============================================================
\section{Fundamentos Cognitivos}
% ============================================================

\begin{frame}{Modelos Mentais na Psicologia}
  \textbf{Kenneth Craik (1943)} --- \textit{The Nature of Explanation}
  \vspace{0.5em}
  \begin{itemize}
    \item Organismos constroem \textbf{``modelos em miniatura''} da realidade
    \item Capacidade de antecipar eventos e testar alternativas mentalmente
    \item Precursor direto dos \textit{world models} em IA
  \end{itemize}
  \vspace{1em}
  \begin{block}{Correlato Neurológico: Vias Visuals}
    \begin{tabularx}{\textwidth}{lX}
      \textbf{Via Ventral} & Reconhecimento de objetos (\textit{``o quê''}) \\
      \textbf{Via Dorsal}  & Processamento espacial (\textit{``onde/como''}) \\
    \end{tabularx}
  \end{block}
  \vspace{0.3em}
  \centering\small Arquiteturas modernas de \textit{world models} frequentemente espelham esta separação.
\end{frame}

\begin{frame}{Princípio da Energia Livre (FEP)}
  \textbf{Karl Friston} --- Sistemas auto-organizáveis minimizam \textit{energia livre variacional} para resistir à entropia.

  \vspace{0.5em}
  Um agente mantém um modelo generativo $p(o_t, s_t)$ sobre observações $o_t$ e estados latentes $s_t$:

  \vspace{0.3em}
  \begin{equation}
    F = \mathbb{E}_{q(s_t)}\bigl[\log q(s_t) - \log p(o_t, s_t)\bigr]
      = D_{\mathrm{KL}}\bigl(q(s_t) \parallel p(s_t \mid o_t)\bigr) - \log p(o_t)
  \end{equation}

  \vspace{0.3em}
  Percepção $\equiv$ inferência variacional: minimizar $F$ maximiza a evidência do modelo interno.
\end{frame}

\begin{frame}{Inferência Ativa}
  Ação e percepção como dois lados da mesma moeda:
  \begin{itemize}
    \item \textbf{Percepção}: atualiza crenças internas $q(s_t)$ para ajustar aos dados $o_t$
    \item \textbf{Ação}: altera o ambiente para forçar dados futuros a se conformarem às expectativas do modelo
  \end{itemize}

  \vspace{0.5em}
  Decomposição da Energia Livre Esperada $G(\pi)$:
  \begin{equation}
    G(\pi) \approx \sum_\tau
      \underbrace{-\mathbb{E}_{q(s_\tau \mid \pi)}\bigl[\mathcal{H}\bigl(q(o_\tau \mid s_\tau)\bigr)\bigr]}_{\text{Epistêmico (Redução de Ambiguidade)}}
      + \underbrace{D_{\mathrm{KL}}\bigl(q(o_\tau \mid \pi) \parallel p(o_\tau)\bigr)}_{\text{Instrumental (Minimização de Risco)}}
  \end{equation}

  \vspace{0.3em}
  Unifica estruturalmente o dilema exploração--explotação do RL \textit{model-free}.
\end{frame}

\begin{frame}{Codificação Preditiva}
  \textbf{Rao \& Ballard (1999)} --- O córtex cerebral implementa uma arquitetura hierárquica que otimiza energia livre via \textit{predictive coding}.

  \vspace{0.3em}
  \begin{center}
    \begin{tikzpicture}[node distance=0.6cm and 0.8cm, auto, >=latex, font=\small]
      \node[layer] (layer2) {Camada Superior $s_{t+1}$};
      \node[layer, below=of layer2, yshift=-0.1cm] (layer1) {Camada Inferior $s_t$};
      \node[error, right=of layer1] (comp) {$\sum$};
      \node[input, below=of comp] (obs) {$o_t$};

      \draw[->, wp-primary, thick] (layer2) -- node[left, font=\tiny] {$f(s_{t+1})$} (layer1);
      \draw[->, wp-primary, thick] (layer1.east) -- (comp.west);
      \draw[->, wp-primary, thick] (obs.north) -- (comp.south);
      \draw[->, wp-primary, thick] (comp.north) .. controls +(up:0.6cm) and +(right:0.6cm) ..
        node[right, midway, font=\tiny] {Erro $\varepsilon_t$} (layer2.east);
    \end{tikzpicture}
  \end{center}

  \vspace{0.3em}
  \begin{itemize}
    \item Predições \textit{top-down} suprimem atividade em camadas inferiores
    \item Apenas erros residuais $\varepsilon_t$ sobem na hierarquia
    \item Mapeia diretamente para RSSMs e PCNs modernos
  \end{itemize}
\end{frame}

% ============================================================
\section{Dois Paradigmas: Preditivo vs Generativo}
% ============================================================

\begin{frame}{Paradigma Preditivo}
  Modela a dinâmida de transição: $p(s_{t+1} \mid s_t, a_t) = f_\theta(s_t, a_t)$

  \vspace{0.3em}
  \begin{block}{Propósito}
    \textbf{Previsão}: dado o que aconteceu, o que acontecerá?
  \end{block}

  \vspace{0.3em}
  \textbf{Representantes:}
  \begin{itemize}
    \item \textbf{Ha \& Schmidhuber (2018)}: VAE + RNN --- arquitetura original \textit{World Models}
    \item \textbf{LeCun (JEPA)}: predição em espaço latente, não em pixels
    \item \textbf{Dreamer (Hafner et al.)}: RSSM + \textit{actor-critic} em trajetórias imaginadas
  \end{itemize}

\end{frame}

\begin{frame}{JEPA --- \textit{Joint Embedding Predictive Architecture}}
  \textbf{LeCun (2022)} --- Estrutura de \textit{self-supervised learning} que prevê em espaço latente:

  \vspace{0.3em}
  \begin{equation}
    \hat{s}_y = \text{Predictor}(s_x, \theta), \quad
    \mathcal{L} = \|\hat{s}_y - s_y\|^2
  \end{equation}

  \begin{itemize}
    \item \textbf{Prevenção de colapso}: \textit{stop-gradient} no encoder alvo
    \item \textbf{Abstração}: descarta detalhes imprevisíveis (ruído de pixel), captura estrutura semântica
    \item \textbf{Jerarquia}: módulos JEPA empilhados operam em diferentes escalas temporais
  \end{itemize}

  \begin{center}
    \begin{tikzpicture}[node distance=0.4cm and 0.8cm, auto, >=latex, font=\small]
      \node[input] (ix) {$x$};
      \node[box, right=of ix] (enc) {Encoder};
      \node[input, right=of enc] (sx) {$s_x$};
      \node[box, right=of sx] (pred) {Predictor};
      \node[input, right=of pred] (syhat) {$\hat{s}_y$};

      \node[input, below=of ix] (iy) {$y$};
      \node[box, right=of iy] (enc2) {Encoder};
      \node[input, right=of enc2] (sg) {$\sg(s_y)$};

      \node[loss, right=of sg] (loss) {$\mathcal{L} = \|\hat{s}_y - s_y\|^2$};

      \draw[->, wp-primary, thick] (ix) -- (enc);
      \draw[->, wp-primary, thick] (enc) -- (sx);
      \draw[->, wp-primary, thick] (sx) -- (pred);
      \draw[->, wp-primary, thick] (pred) -- (syhat);
      \draw[->, wp-primary, thick] (iy) -- (enc2);
      \draw[->, wp-primary, thick] (enc2) -- (sg);
      \draw[->, wp-primary, thick] (syhat.south) -| (loss.north);
      \draw[->, wp-primary, thick] (sg) -- (loss);
      \draw[wp-muted, dashed, thick] (enc.south) -- (enc2.north)
        node[midway, right, wp-muted, font=\tiny] {pesos compartilhados};
    \end{tikzpicture}
  \end{center}
\end{frame}

\begin{frame}{Paradigma Generativo}
  Modela a distribuição conjunta completa:
  \begin{equation}
    p(o_1, \dots, o_T, a_1, \dots, a_T) = \prod_{t=1}^T p(o_t, a_t \mid o_{<t}, a_{<t})
  \end{equation}

  \vspace{0.3em}
  \textbf{Representantes:}
  \begin{itemize}
    \item \textbf{Google Genie}: tokenizador de vídeo + modelo de dinâmica causal; mundo jogável a partir de uma única imagem
    \item \textbf{GAIA-1 / UniSim}: transformers de vídeo em larga escala que aprendem física intuitiva diretamente de pixels
  \end{itemize}

  \vspace{0.5em}
  \begin{block}{Avaliação}
    Qualidade da amostra, diversidade distribucional, consistência temporal de longo horizonte
  \end{block}
\end{frame}

\begin{frame}{Paradigmas em Contraste}
  \centering
  \footnotesize
  \begin{tabular}{lp{4.2cm}p{4.2cm}}
    \toprule
    \textbf{Dimensão} & \textbf{Preditivo} & \textbf{Generativo} \\
    \midrule
    Objetivo Primário & Minimizar distância latente & Maximizar log-verossimilhança \\
    Tipo de Saída & Predição de estado abstrato & Amostragem do fluxo observacional \\
    Sinal de Aprendizagem & Alinhamento de atributos & Reconstrução / \textit{token loss} \\
    Utilidade & Planejamento \textit{local look-ahead} & Simulação ambiental \\
    Vulnerabilidade & Colapso representacional & Alto custo amostral \\
    \bottomrule
  \end{tabular}

  \vspace{0.5em}
  \pause
  \begin{alertblock}{}
    A fronteira é fluida --- Dreamer usa um modelo generativo para \textit{rollouts} preditivos; Genie integra um controlador preditivo em um pipeline generativo.
  \end{alertblock}
\end{frame}

% ============================================================
\section{Fundações Teóricas de Schmidhuber}
% ============================================================

\begin{frame}{Schmidhuber: Bases Teóricas}
  \vspace{-0.3em}
  \begin{block}{Curiosidade Artificial (1991)}
    Recompensa intrínseca por erros de predição \textbf{aprendíveis}:
    \begin{equation}
      \text{curiosity reward} = \|\hat{s}_{t+1} - s_{t+1}\|^2 \quad (\text{apenas se aprendizável})
    \end{equation}
    Inspirou módulos de curiosidade intrínseca em RL moderno (ICM, Pathak et al. 2017).
  \end{block}

  \begin{block}{Teoria Formal da Criatividade (2012)}
    \begin{itemize}
      \item ``Beleza'' $\propto$ progresso de compressão/aprendizagem no \textit{world model} do observador
      \item Obras de arte, ciência e humor geram progresso de compressão aprendizável
    \end{itemize}
  \end{block}
\end{frame}

% ============================================================
\section{Arquiteturas Fundamentais}
% ============================================================

\begin{frame}{Ha \& Schmidhuber (2018): Arquitetura VMC}
  \begin{center}
    \begin{tikzpicture}[node distance=0.6cm, auto, >=latex]
      \node[comp] (v) {Visão (V) \\ \footnotesize VAE};
      \node[comp, below=of v] (m) {Memória (M) \\ \footnotesize RNN};
      \node[comp, below=of m] (c) {Controlador (C) \\ \footnotesize Política};

      \draw[->, wp-primary, thick] (v) -- (m);
      \draw[->, wp-primary, thick] (m) -- (c);
      \draw[->, wp-primary, thick] (c.east) .. controls +(right:0.8cm) and +(right:0.8cm) .. (v.east);
    \end{tikzpicture}
  \end{center}

  \begin{enumerate}
    \item \textbf{Visão (V)}: VAE comprime observações $x_t$ em latente $z_t \sim \mathcal{N}(\mu_t, \sigma_t)$
    \item \textbf{Memória (M)}: RNN prediz próximo estado latente $\hat{z}_{t+1} = f_\theta(z_t, a_t)$
    \item \textbf{Controlador (C)}: Política $a_t = \pi(z_t)$ treinada no espaço latente
  \end{enumerate}

  \vspace{0.1em}
  Planejamento por \textit{``sonhar''} --- \textit{rollouts} no espaço latente sem simulação em nível de pixel.
\end{frame}

\begin{frame}{RSSM --- \textit{Recurrent State-Space Model}}
  \textbf{Hafner et al. (Dreamer)} --- Combina estados determinísticos e estocásticos:

  \vspace{0.3em}
  \begin{align}
    h_t &= f_\phi(h_{t-1}, z_{t-1}, a_{t-1}) &&\text{(determinístico)} \\
    z_t &\sim \mathcal{N}\bigl(\mu_\phi(h_t), \sigma_\phi(h_t)\bigr) &&\text{(estocástico)} \\
    \hat{o}_t &= g_\phi(h_t, z_t) &&\text{(reconstrução)}
  \end{align}

  \vspace{0.3em}
  \begin{itemize}
    \item $h_t$ carrega dependências temporais de longo alcance
    \item $z_t$ captura futuros multimodais
    \item Evita colapso posterior de modelos puramente estocásticos
  \end{itemize}
\end{frame}

\begin{frame}{DreamerV2 e V3}
  \begin{columns}
    \begin{column}{0.48\textwidth}
      \textbf{DreamerV2} (Hafner et al., 2021)
      \begin{itemize}
        \item Latentes \textbf{discretos} (categoriais) ao invés de Gaussianos
        \item \textbf{KL balancing}: coeficientes separados para \textit{prior} e \textit{posterior}
        \item Predição de recompensa \textit{two-hot}
      \end{itemize}
      \begin{equation}
        z_t = \text{one-hot}(\argmax_k \pi_k) + \pi - \sg(\pi)
      \end{equation}
    \end{column}
    \begin{column}{0.48\textwidth}
      \textbf{DreamerV3} (Hafner et al., 2023)
      \begin{itemize}
        \item Hiperparâmetros únicos para centenas de ambientes
        \item \textbf{Predições Symlog}: $\symlog(x) = \text{sign}(x)\log(|x|+1)$
        \item \textbf{Free bits}: regularizador KL que evita colapso
        \item Predições normalizadas por estatísticas correntes
      \end{itemize}
    \end{column}
  \end{columns}
\end{frame}

% ============================================================
\section{World Models Baseados em Transformers}
% ============================================================

\begin{frame}{IRIS --- \textit{Transformers are Sample Efficient World Models}}
  \textbf{Micheli et al. (2023)} --- Substitui memória recorrente por transformers:

  \vspace{0.5em}
  \begin{center}
    \begin{tikzpicture}[node distance=0.3cm and 0.6cm, auto, >=latex, font=\small]
      \node[input] (obs) {$o_t$};
      \node[box, right=of obs] (vq) {VQ-VAE};
      \node[input, right=of vq] (tok) {$z_t$};
      \node[box, right=of tok] (trans) {Transformer};
      \node[input, right=of trans] (pred) {$\hat{z}_{t+1}$};
      \node[box, right=of pred] (dec) {Decode};
      \node[input, right=of dec] (out) {$\hat{o}_{t+1}$};

      \draw[->, wp-primary, thick] (obs) -- (vq);
      \draw[->, wp-primary, thick] (vq) -- (tok);
      \draw[->, wp-primary, thick] (tok) -- (trans);
      \draw[->, wp-primary, thick] (trans) -- (pred);
      \draw[->, wp-primary, thick] (pred) -- (dec);
      \draw[->, wp-primary, thick] (dec) -- (out);
    \end{tikzpicture}
  \end{center}

  \vspace{0.3em}
  \begin{equation}
    p(z_{t+1} \mid z_{\leq t}, a_{\leq t}) = \prod_{k=1}^K p(z_{t+1}^k \mid z_{< t+1}^k, z_{\leq t}, a_{\leq t})
  \end{equation}

  \vspace{0.3em}
  \begin{itemize}
    \item Observação $o_t$ tokenizada por VQ-VAE em códigos discretos $z_t$
    \item Transformer autoregressivo prediz a próxima sequência de tokens
    \item Eficiência amostral competitiva com métodos \textit{model-free}
  \end{itemize}
\end{frame}

% ============================================================
\section{Aplicações}
% ============================================================

\begin{frame}{Aplicações de \textit{World Models}}
  \begin{columns}
    \begin{column}{0.32\textwidth}
      \textbf{Robótica}
      \begin{itemize}\small
        \item \textit{DayDreamer}: DreamerV3 em robôs reais
        \item Políticas treinadas na ``imaginação''
        \item Mínima interação com o mundo físico
      \end{itemize}
    \end{column}
    \begin{column}{0.32\textwidth}
      \textbf{Direção Autônoma}
      \begin{itemize}\small
        \item Wayve: simulação de direção
        \item Avaliação \textit{closed-loop} de políticas
        \item Segurança: teste de cenários críticos sem risco real
      \end{itemize}
    \end{column}
    \begin{column}{0.32\textwidth}
      \textbf{RL Baseado em Modelo}
      \begin{itemize}\small
        \item \textit{Rollouts} alucinados
        \item Eficiência amostral: ordens de grandeza
        \item Planejamento ciente de incerteza (ensembles)
      \end{itemize}
    \end{column}
  \end{columns}
\end{frame}

% ============================================================
\section{Métricas de Avaliação}
% ============================================================

\begin{frame}{Métricas de Avaliação}
  \centering
  \begin{tabularx}{\textwidth}{lX}
    \toprule
    \textbf{Métrica} & \textbf{Descrição} \\
    \midrule
    \textbf{Prediction MSE} & Erro quadrático médio em nível de pixel ou latente \\
    \textbf{Divergência KL} & Casamento entre \textit{prior} e \textit{posterior} latentes \\
    \textbf{Acurácia de Planejamento} & Correlação entre trajetórias imaginadas e retornos reais \\
    \textbf{Eficiência Amostral} & Passos de ambiente necessários para atingir \textit{threshold} de recompensa \\
    \textbf{Gap de Generalização} & Queda de performance em variações não vistas do ambiente \\
    \bottomrule
  \end{tabularx}
\end{frame}

% ============================================================
\section{Problemas em Aberto}
% ============================================================

\begin{frame}{Problemas em Aberto}
  \begin{description}
    \item[Composicionalidade] Como construir \textit{world models} que combinam conceitos conhecidos para generalizar a cenários novos?
    \item[Predição de Longo Horizonte] Erros de predição se acumulam quadraticamente --- difusão iterativa, atenção esparsa e separação hierárquica de escalas são soluções exploradas.
    \item[Escalabilidade] Ambientes controlados (ViZDoom, CarRacing) $\rightarrow$ mundo real aberto com observações de alta dimensão.
    \item[Grounding] Como manter significado consistente das representações latentes entre diferentes tarefas?
    \item[Causalidade] Distinguir intervenção de observação para raciocínio contrafactual genuíno.
  \end{description}
\end{frame}

% ============================================================
\section{Conclusão}
% ============================================================

\begin{frame}{Conclusão}
  \vspace{-0.8em}
  \begin{block}{Trajetória}
    \begin{center}
      \begin{tikzpicture}[node distance=0.25cm, auto, >=latex, font=\tiny]
        \node[era, minimum width=1.6cm, minimum height=0.35cm, font=\tiny] (craik) {1943: Craik};
        \node[era, minimum width=1.6cm, minimum height=0.35cm, font=\tiny, below=of craik] (pred) {1999: Pred. Coding};
        \node[era, minimum width=1.6cm, minimum height=0.35cm, font=\tiny, below=of pred] (wm) {2018: W. Models};
        \node[era, minimum width=1.6cm, minimum height=0.35cm, font=\tiny, below=of wm] (dream) {2020: Dreamer};
        \node[era, minimum width=1.6cm, minimum height=0.35cm, font=\tiny, below=of dream] (v2) {2021: DreamerV2};
        \node[era, minimum width=1.6cm, minimum height=0.35cm, font=\tiny, below=of v2] (iris) {2023: IRIS / V3};

        \draw[->, wp-muted, thick] (craik) -- (pred);
        \draw[->, wp-muted, thick] (pred) -- (wm);
        \draw[->, wp-primary, thick] (wm) -- (dream);
        \draw[->, wp-primary, thick] (dream) -- (v2);
        \draw[->, wp-primary, thick] (v2) -- (iris);
      \end{tikzpicture}
    \end{center}
  \end{block}

  \vspace{-0.5em}
  \begin{itemize}\small
    \item \textit{World models} representam a convergência de ciência cognitiva, teoria de controle e \textit{deep learning}
    \item Desafios-chave: composicionalidade, longo horizonte, escalabilidade, causalidade
    \item Papel central na próxima geração de sistemas inteligentes
  \end{itemize}
\end{frame}

\end{document}
